% =====================================================================
%  Chapitre 5 : Balises et filtres
% =====================================================================
\chapter{Les balises et les filtres}
\label{ch:balises}

\chapitreresume{Le langage de gabarit de Django est volontairement limité. Il
fournit des balises pour les boucles, les conditions et les URL, et des filtres
pour transformer les valeurs à l'affichage. Dans ce chapitre, vous apprenez les
balises et filtres essentiels, puis vous créez vos propres filtres et balises
réutilisables.}

\begin{objectifs}
\begin{itemize}
  \item Distinguer une variable, une balise et un filtre.
  \item Utiliser les balises \code{if}, \code{for}, \code{with}, \code{url},
        \code{include} et \code{csrf\_token}.
  \item Appliquer les filtres courants : \code{date}, \code{length},
        \code{truncatewords}, \code{pluralize}, \code{yesno}.
  \item Créer un filtre personnalisé et une balise simple.
  \item Créer une balise d'inclusion pour réutiliser une carte de livre.
  \item Comprendre l'échappement automatique et éviter les failles XSS.
\end{itemize}
\end{objectifs}

\section{La syntaxe du langage de gabarit}

Le langage de gabarit de Django utilise quatre formes :

\begin{codebox}[html]{Les quatre formes de base}
{# Ceci est un commentaire court #}

{{ livre.titre }}

{{ livre.titre|upper }}

{{ livre.resume|truncatewords:20 }}

{% if livre.est_disponible %}
  <p>Disponible</p>
{% endif %}
\end{codebox}

\begin{description}[leftmargin=3.2cm,style=nextline]
  \item[\code{\{\{ ... \}\}}] Affiche une valeur. On parle de variable ou
  d'expression.
  \item[\code{\{\% ... \%\}}] Exécute une balise : boucle, condition, URL,
  inclusion, chargement de bibliothèque.
  \item[\code{\{\# ... \#\}}] Commentaire court, sur une seule ligne. Il n'est pas
  envoyé au navigateur.
  \item[\code{|filtre}] Transforme la valeur affichée. Les filtres se chaînent :
  \code{nom|lower|capfirst}.
\end{description}

\begin{retenir}
Le langage de gabarit n'est pas Python. Il ne permet pas d'appeler des fonctions
avec des arguments, d'importer des modules ou d'écrire des boucles complexes.
Cette limitation est volontaire : elle protège la séparation entre logique et
présentation.
\end{retenir}

\section{Les variables et la recherche par point}

Dans \code{\{\{ livre.auteur.nom \}\}}, Django cherche dans cet ordre :
\begin{enumerate}
  \item la clé \code{"auteur"} dans un dictionnaire ;
  \item l'attribut \code{auteur} d'un objet ;
  \item la méthode \code{auteur()} sans argument ;
  \item l'index \code{auteur} d'une liste.
\end{enumerate}

Puis il répète l'opération pour \code{.nom}. Cette recherche par point explique
pourquoi \code{\{\{ livre.titre \}\}} fonctionne aussi bien avec un dictionnaire
qu'avec un objet \code{Livre}.

\begin{codebox}[html]{Exemples de recherche}
{{ livre.titre }}
{{ livre.auteur.nom }}
{{ livre.categories.all }}
{{ livres.0.titre }}
{{ user.get_full_name }}
\end{codebox}

\begin{attention}
Si une variable n'existe pas, Django affiche une chaîne vide, sans erreur. C'est
pratique, mais cela peut masquer une faute de frappe. Si la page est vide à
l'endroit attendu, vérifiez le nom exact de la variable dans la vue.
\end{attention}

\begin{note}
\code{\{\{ user.get\_full\_name \}\}} appelle la méthode \code{get\_full\_name()} sans
parenthèses. Le langage de gabarit ajoute automatiquement les parenthèses pour les
méthodes qui n'attendent aucun argument. Une méthode qui exige un argument ne peut
pas être appelée directement.
\end{note}

\section{Les balises essentielles}

\subsection{if, elif, else}

\begin{codebox}[html]{Conditions}
{% if livre.est_disponible %}
  <span class="badge badge-disponible">Disponible</span>
{% elif livre.statut == "reserve" %}
  <span class="badge badge-indisponible">Réservé</span>
{% else %}
  <span class="badge badge-indisponible">Indisponible</span>
{% endif %}
\end{codebox}

Les opérateurs disponibles sont \code{==}, \code{!=}, \code{<}, \code{>},
\code{<=}, \code{>=}, \code{in}, \code{not in}, \code{is}, \code{is not},
\code{and}, \code{or} et \code{not}. Les parenthèses ne sont pas autorisées : on
combine avec \code{and} et \code{or}, en gardant des conditions simples.

\begin{attention}
\code{\{\% if livre \}\}} teste la vérité de l'objet. Un objet Django est
généralement vrai, même s'il n'a pas de titre. Pour tester une chaîne vide,
utilisez \code{\{\% if livre.titre \}\}}. Pour tester une liste vide, utilisez
\code{\{\% if livres \}\}}.
\end{attention}

\subsection{for et empty}

\begin{codebox}[html]{Boucle avec forloop}
<ul>
  {% for livre in livres %}
    <li>
      {{ forloop.counter }}. {{ livre.titre }}
      {% if forloop.first %}(nouveauté){% endif %}
      {% if forloop.last %}(dernier de la liste){% endif %}
    </li>
  {% empty %}
    <li>Aucun livre pour le moment.</li>
  {% endfor %}
</ul>
\end{codebox}

\begin{description}[leftmargin=3.6cm,style=nextline]
  \item[\code{forloop.counter}] Compteur à partir de 1.
  \item[\code{forloop.counter0}] Compteur à partir de 0.
  \item[\code{forloop.revcounter}] Compteur à rebours, fin à 1.
  \item[\code{forloop.first}] Vrai au premier tour.
  \item[\code{forloop.last}] Vrai au dernier tour.
  \item[\code{forloop.parentloop}] Boucle parente, dans une boucle imbriquée.
\end{description}

On peut parcourir un dictionnaire :

\begin{codebox}[html]{Boucle sur un dictionnaire}
{% for cle, valeur in statistiques.items %}
  <p>{{ cle }} : {{ valeur }}</p>
{% endfor %}
\end{codebox}

Et inverser une liste avec le filtre \code{reversed} :

\begin{codebox}[html]{Boucle inversée}
{% for livre in livres reversed %}
  <p>{{ livre.titre }}</p>
{% endfor %}
\end{codebox}

\subsection{with}

\code{with} crée un alias pour une expression longue ou coûteuse :

\begin{codebox}[html]{Alias avec with}
{% with auteur=livre.auteur.nom_complet %}
  <p>Écrit par {{ auteur }}.</p>
  <p>Initiales : {{ auteur|initiales }}.</p>
{% endwith %}
\end{codebox}

\subsection{url}

\begin{codebox}[html]{Construire une URL}
<a href="{% url 'catalogue:accueil' %}">Accueil</a>
<a href="{% url 'catalogue:livre_detail' livre.slug %}">Voir le livre</a>
\end{codebox}

La balise \code{url} accepte aussi des paramètres nommés :

\begin{codebox}[html]{Paramètres nommés}
<a href="{% url 'catalogue:livre_detail' slug=livre.slug %}">
  Voir le livre
</a>
\end{codebox}

\subsection{csrf\_token}

Tout formulaire \code{POST} doit contenir \lstinline|{% csrf_token %}|. Sans ce
jeton, Django refuse la requête avec une erreur \code{403 Forbidden}.

\begin{codebox}[html]{Jeton CSRF}
<form method="post">
  {% csrf_token %}
  <input type="text" name="titre">
  <button type="submit">Enregistrer</button>
</form>
\end{codebox}

\begin{attention}
\code{csrf\_token} protège contre les attaques CSRF, qui consistent à faire
envoyer un formulaire à un utilisateur connecté depuis un autre site. Ne
désactivez jamais cette protection sans raison très précise.
\end{attention}

\subsection{include, extends et block}

\begin{codebox}[html]{Inclusion et héritage}
{% extends "base.html" %}

{% block contenu %}
  {% include "partials/entete.html" %}
  <p>Contenu de la page.</p>
{% endblock %}
\end{codebox}

\code{include} insère un gabarit complet. \code{extends} établit une relation
d'héritage. Un gabarit enfant ne peut pas utiliser \code{extends} et être inclus
en même temps : choisissez l'un ou l'autre.

\subsection{comment, autoescape, now, firstof, cycle}

\begin{codebox}[html]{Balises utilitaires}
{% comment %}
  Ce bloc est ignoré par Django.
  Utile pour désactiver temporairement du code.
{% endcomment %}

{% autoescape off %}
  {{ contenu_html }}
{% endautoescape %}

<p>Nous sommes le {% now "d/m/Y" %}.</p>

<p>{{ livre.sous_titre|default:livre.titre }}</p>

{% firstof livre.sous_titre livre.titre "Sans titre" %}

{% for i in "123" %}
  <span class="{% cycle 'pair' 'impair' %}">{{ i }}</span>
{% endfor %}
\end{codebox}

\begin{attention}
\code{autoescape off} désactive l'échappement HTML pour toute la zone. Ne
l'utilisez jamais avec du contenu fourni par un utilisateur. Nous y revenons à la
fin du chapitre.
\end{attention}

\subsection{verbatim}

Si vous devez afficher des accolades qui ressemblent à des balises Django, par
exemple dans du JavaScript ou un exemple de code, utilisez \code{verbatim} :

\begin{codebox}[html]{Afficher des accolades}
{% verbatim %}
  <p>{{ ceci_est_affiche_tel_quel }}</p>
{% endverbatim %}
\end{codebox}

\section{Les filtres essentiels}

Un filtre se place après une valeur, séparé par une barre verticale. Certains
acceptent un argument après deux-points.

\begin{center}
\footnotesize
\begin{tabularx}{\textwidth}{L{8.8cm} X}
\toprule
\textbf{Expression} & \textbf{Résultat} \\
\midrule
\code{\{\{ livre.titre|upper \}\}} & Le titre en majuscules. \\
\code{\{\{ livre.titre|lower \}\}} & Le titre en minuscules. \\
\code{\{\{ livre.titre|title \}\}} & Chaque mot commence par une majuscule. \\
\code{\{\{ livre.titre|capfirst \}\}} & Seule la première lettre est en majuscule. \\
\code{\{\{ livre.titre|length \}\}} & Le nombre de caractères. \\
\code{\{\{ livres|length \}\}} & Le nombre d'éléments de la liste. \\
\code{\{\{ livre.resume|truncatechars:80 \}\}} & Coupe à 80 caractères. \\
\code{\{\{ livre.resume|truncatewords:15 \}\}} & Coupe à 15 mots. \\
\code{\{\{ livre.date\_publication|date:"d/m/Y" \}\}} & Une date au format français. \\
\code{\{\{ livre.prix|floatformat:2 \}\}} & Deux chiffres après la virgule. \\
\code{\{\{ livre.categories.all|join:", " \}\}} & Joint les catégories par une virgule. \\
\code{\{\{ livre.nombre\_exemplaires|pluralize \}\}} & Ajoute \og s \fg{} si le nombre est différent de 1. \\
\code{\{\{ livre.est\_disponible|yesno:"Oui,Non" \}\}} & Affiche Oui ou Non. \\
\code{\{\{ livre.titre|slugify \}\}} & Transforme le titre en slug. \\
\code{\{\{ livre.resume|linebreaks \}\}} & Transforme les sauts de ligne en paragraphes. \\
\code{\{\{ livre.resume|striptags \}\}} & Supprime les balises HTML. \\
\bottomrule
\end{tabularx}
\end{center}

\subsection{default et default\_if\_none}

\begin{codebox}[html]{Valeurs par défaut}
{{ livre.sous_titre|default:"Sans sous-titre" }}
{{ livre.editeur|default_if_none:"Éditeur inconnu" }}
{{ livre.annee|default:"Année inconnue" }}
\end{codebox}

\code{default} remplace les valeurs fausses : chaîne vide, zéro, liste vide,
\code{None}. \code{default\_if\_none} ne remplace que \code{None}. La différence
compte pour un nombre : \code{0|default:"Non renseigné"} affiche
\og Non renseigné \fg, alors que \code{0|default\_if\_none:"Non renseigné"}
affiche \code{0}.

\subsection{date et time}

\begin{codebox}[html]{Formats de date}
{{ livre.date_publication|date:"d/m/Y" }}
{{ livre.date_publication|date:"l j F Y" }}
{{ emprunt.date_emprunt|date:"D d M Y" }}
{{ maintenant|time:"H:i" }}
\end{codebox}

Les formats les plus utilisés : \code{d} (jour sur deux chiffres), \code{j}
(jour sans zéro), \code{m} (mois sur deux chiffres), \code{F} (mois en toutes
lettres), \code{Y} (année sur quatre chiffres), \code{H} (heure sur 24 heures),
\code{i} (minutes), \code{s} (secondes).

\subsection{pluralize et yesno}

\begin{codebox}[html]{Accords et booléens}
<p>
  {{ page_obj.paginator.count }} livre{{ page_obj.paginator.count|pluralize }}
</p>
<p>{{ livre.est_disponible|yesno:"Disponible,Indisponible" }}</p>
\end{codebox}

\code{pluralize} ajoute un \og s \fg{} par défaut. On peut préciser deux formes :
\code{pluralize:"y,ies"} pour \og copy \fg{} et \og copies \fg.

\subsection{linebreaks, escape et safe}

\begin{codebox}[html]{Texte et HTML}
{{ livre.resume|linebreaks }}
{{ livre.resume|linebreaksbr }}
{{ commentaire|escape }}
{{ contenu|safe }}
\end{codebox}

\begin{attention}
\code{linebreaks} échappe le contenu, puis ajoute des balises \code{<p>} et
\code{<br>}. \code{safe} ne fait aucun échappement : il déclare que la chaîne est
déjà du HTML sûr. Utilisez \code{safe} uniquement sur du contenu que vous
contrôlez entièrement.
\end{attention}

\section{Créer ses propres filtres et balises}

Les filtres intégrés couvrent la majorité des besoins. Pour les cas spécifiques,
Django permet d'écrire ses propres bibliothèques de balises.

\subsection{L'arborescence}

Créez un dossier \code{templatetags/} dans l'application, avec un fichier
\code{\_\_init\_\_.py} vide :

\begin{codebox}[tree]{Arborescence des balises personnalisées}
catalogue/
|-- templatetags/
|   |-- __init__.py
|   `-- biblio.py
|-- models.py
`-- views.py
\end{codebox}

\begin{attention}
Le dossier doit s'appeler exactement \code{templatetags}, au pluriel, et se trouver
dans l'application. Le fichier \code{\_\_init\_\_.py} est obligatoire : sans lui,
Django ne voit pas le paquet. Redémarrez le serveur après la création.
\end{attention}

\subsection{Un filtre personnalisé}

\begin{codebox}[python]{catalogue/templatetags/biblio.py}
from django import template

register = template.Library()


@register.filter
def initiales(nom_complet):
    """Retourne les initiales d'un nom, par exemple AS."""
    morceaux = nom_complet.split()
    return "".join(morceau[0] for morceau in morceaux[:2]).upper()
\end{codebox}

Utilisation dans un gabarit :

\begin{codebox}[html]{Utilisation du filtre}
{% load biblio %}
<p>{{ livre.auteur }} ({{ livre.auteur|initiales }})</p>
\end{codebox}

\begin{note}
Le nom du filtre est le nom de la fonction. On peut le changer avec
\code{@register.filter(name="sigle")}. Le nom doit être un identifiant valide et
ne pas contenir d'espace.
\end{note}

\subsection{Une balise simple}

Une balise simple renvoie une valeur. Elle peut recevoir des arguments.

\begin{codebox}[python]{catalogue/templatetags/biblio.py, extrait}
@register.simple_tag
def cote(auteur, annee):
    """Construit une cote simple, par exemple S1943."""
    return f"{auteur[0].upper()}{annee}"
\end{codebox}

\begin{codebox}[html]{Utilisation de la balise}
{% load biblio %}
<p>Cote : {% cote livre.auteur.nom livre.annee %}</p>
\end{codebox}

\subsection{Une balise d'inclusion}

Une balise d'inclusion rend un gabarit avec un contexte. C'est la manière propre
de réutiliser un composant, comme une carte de livre.

\begin{codebox}[python]{catalogue/templatetags/biblio.py, extrait}
@register.inclusion_tag("catalogue/livres/carte.html")
def carte_livre(livre):
    return {"livre": livre}
\end{codebox}

\begin{codebox}[html]{templates/catalogue/livres/carte.html}
{% load biblio %}
<article class="carte-livre">
  <div class="couverture-vide" aria-hidden="true">
    {{ livre.titre|first }}
  </div>
  <h3>{{ livre.titre|truncatechars:38 }}</h3>
  <p class="auteur">{{ livre.auteur }}</p>
  <span class="badge badge-disponible">
    {{ livre.annee|default:"Année inconnue" }}
  </span>
</article>
\end{codebox}

\begin{codebox}[html]{Utilisation dans la page d'accueil}
{% load biblio %}

<div class="grille-livres">
  {% for livre in livres %}
    {% carte_livre livre %}
  {% endfor %}
</div>
\end{codebox}

\begin{description}[leftmargin=3.5cm,style=nextline]
  \item[\code{inclusion\_tag}] Déclare une balise qui rend un gabarit. Le
  dictionnaire renvoyé devient le contexte du gabarit.
  \item[\code{takes\_context=True}] Permet à la balise de recevoir le contexte
  complet, par exemple pour lire \code{request.GET}.
\end{description}

\subsection{Une balise avec contexte}

Pour construire un lien de pagination qui conserve les filtres de recherche, la
balise a besoin de la requête :

\begin{codebox}[python]{catalogue/templatetags/biblio.py, extrait}
@register.simple_tag(takes_context=True)
def url_page(context, numero):
    """Construit un lien ?page=N en conservant les autres paramètres."""
    params = context["request"].GET.copy()
    params["page"] = numero
    return f"?{params.urlencode()}"
\end{codebox}

\begin{codebox}[html]{Utilisation}
<a href="{% url_page page_obj.previous_page_number %}">Précédent</a>
<a href="{% url_page page_obj.next_page_number %}">Suivant</a>
\end{codebox}

\begin{note}
\code{context["request"]} est disponible grâce au processeur de contexte
\code{django.template.context\_processors.request}, présent dans
\code{settings.py}. Sans lui, la balise lèverait une erreur
\code{KeyError: 'request'}.
\end{note}

\section{Intégrer les composants au projet}

Mettez à jour la vue d'accueil pour enrichir les données :

\begin{codebox}[python]{catalogue/views.py, extrait}
LIVRES_DEMO = [
    {
        "titre": "Le Petit Prince",
        "auteur": "Antoine de Saint-Exupéry",
        "annee": 1943,
        "resume": "Un aviateur rencontre un petit prince venu d'une autre planète.",
    },
    {
        "titre": "Une si longue lettre",
        "auteur": "Mariama Bâ",
        "annee": 1979,
        "resume": "Une lettre qui raconte la vie de femmes sénégalaises.",
    },
    {
        "titre": "L'Étranger",
        "auteur": "Albert Camus",
        "annee": 1942,
        "resume": "Meursault raconte un absurde fait divers sous le soleil d'Alger.",
    },
]


def accueil(request):
    return render(
        request,
        "catalogue/accueil.html",
        {"livres": LIVRES_DEMO},
    )


def livre_liste(request):
    return render(
        request,
        "catalogue/livres/liste.html",
        {"livres": LIVRES_DEMO},
    )
\end{codebox}

Puis le gabarit :

\begin{codebox}[html]{templates/catalogue/accueil.html}
{% extends "base.html" %}
{% load biblio %}

{% block titre %}Accueil | GesiBiblio{% endblock %}

{% block contenu %}
  <section class="hero">
    <div>
      <p class="surtitre">Bibliothèque numérique</p>
      <h1>Empruntez, lisez, revenez.</h1>
      <p>
        {{ livres|length }} livre{{ livres|length|pluralize }} vous attendent
        dans le catalogue.
      </p>
      <a class="bouton bouton-principal"
         href="{% url 'catalogue:livre_liste' %}">Explorer le catalogue</a>
    </div>
  </section>

  <h2>Nouveautés</h2>
  <div class="grille-livres">
    {% for livre in livres %}
      {% carte_livre livre %}
    {% empty %}
      <p>Aucun livre pour le moment.</p>
    {% endfor %}
  </div>
{% endblock %}
\end{codebox}

Ajoutez la route du catalogue dans \code{catalogue/urls.py} :

\begin{codebox}[python]{catalogue/urls.py, extrait}
path("livres/", views.livre_liste, name="livre_liste"),
\end{codebox}

Puis créez la page de liste. Elle réutilise la balise \code{carte\_livre} et
montre comment parcourir une collection :

\begin{codebox}[html]{templates/catalogue/livres/liste.html}
{% extends "base.html" %}
{% load biblio %}

{% block titre %}Catalogue | GesiBiblio{% endblock %}

{% block contenu %}
  <h1>Catalogue</h1>
  <p>
    {{ livres|length }} livre{{ livres|length|pluralize }}
    dans le catalogue.
  </p>

  <div class="grille-livres">
    {% for livre in livres %}
      {% carte_livre livre %}
    {% empty %}
      <p>Aucun livre ne correspond à votre recherche.</p>
    {% endfor %}
  </div>
{% endblock %}
\end{codebox}

\begin{note}
La recherche, les filtres et la pagination seront ajoutés au chapitre de la
gestion des requêtes. Pour l'instant, la page affiche toute la liste de
démonstration.
\end{note}

Ajoutez un petit style pour les initiales si vous les affichez :

\begin{codebox}[css]{static/css/style.css, extrait}
.avatar {
  display: inline-grid;
  place-items: center;
  width: 2rem;
  height: 2rem;
  border-radius: 50%;
  background: var(--navy);
  color: var(--white);
  font-size: 0.72rem;
  font-weight: 800;
}
\end{codebox}

\section{Sécurité : échappement et XSS}

Django échappe automatiquement les variables. Si un utilisateur enregistre une
balise \code{<script>} contenant du code JavaScript dans le résumé d'un livre, le
gabarit affiche le texte au lieu de l'exécuter.

\begin{codebox}[html]{Échappement automatique}
{{ livre.resume }}
\end{codebox}

Le navigateur reçoit :

\begin{codebox}[html]{HTML produit}
&lt;script&gt;alert("xss")&lt;/script&gt;
\end{codebox}

Avec \code{|safe}, le script serait exécuté. C'est une faille XSS.

\begin{codebox}[html]{À ne jamais faire avec une donnée utilisateur}
{{ livre.resume|safe }}
\end{codebox}

\begin{attention}
Les filtres \code{escape} et l'échappement automatique sont vos amis. Si vous
avez besoin d'afficher du HTML riche, utilisez une bibliothèque de nettoyage comme
\code{bleach} ou \code{nh3}, puis marquez le résultat comme sûr uniquement après
nettoyage. Ne faites jamais confiance à une chaîne envoyée par un utilisateur.
\end{attention}

\begin{bonne}
Règle simple : pas de \code{safe} sur les champs remplis par les utilisateurs. Pour
les sauts de ligne, utilisez \code{linebreaks} ou \code{linebreaksbr}, qui
échappent le contenu avant d'ajouter les balises.
\end{bonne}

\section{Exercices}

\begin{exercice}[ : compter et accorder]
Affichez le nombre de livres de la page d'accueil avec \code{length} et
\code{pluralize}. Le texte doit être correct pour 0, 1 et plusieurs livres.
\end{exercice}

\begin{correction}
\begin{codebox}[html]{Extrait}
<p>
  {{ livres|length }}
  livre{{ livres|length|pluralize }}
  dans la sélection.
</p>
\end{codebox}

Avec 0 livre, \code{pluralize} renvoie \og s \fg{} car 0 est différent de 1. Le
texte affiche donc \og 0 livres \fg. Avec 1, il affiche \og 1 livre \fg.
\end{correction}

\begin{exercice}[ : filtre extrait]
Créez un filtre \code{extrait} qui renvoie les 90 premiers caractères d'un texte,
suivis de points de suspension si le texte est plus long.
\end{exercice}

\begin{correction}
\begin{codebox}[python]{catalogue/templatetags/biblio.py, extrait}
@register.filter
def extrait(texte, longueur=90):
    texte = str(texte)
    if len(texte) <= longueur:
        return texte
    return texte[:longueur].rstrip() + "..."
\end{codebox}

\begin{codebox}[html]{Utilisation}
{{ livre.resume|extrait:70 }}
\end{codebox}

Le filtre accepte un argument optionnel. \code{rstrip()} évite de couper au milieu
d'un espace.
\end{correction}

\begin{exercice}[ : balise de date]
Créez une balise simple \code{aujourdhui} qui affiche la date du jour au format
\code{d/m/Y}. Utilisez \code{django.utils.timezone.localdate}.
\end{exercice}

\begin{correction}
\begin{codebox}[python]{catalogue/templatetags/biblio.py, extrait}
from django.utils import timezone


@register.simple_tag
def aujourdhui():
    return timezone.localdate().strftime("%d/%m/%Y")
\end{codebox}

\begin{codebox}[html]{Utilisation}
{% load biblio %}
<p>Mis à jour le {% aujourdhui %}.</p>
\end{codebox}
\end{correction}

\begin{retenir}
\begin{itemize}
  \item \code{\{\{ ... \}\}} affiche, \code{\{\% ... \%\}} exécute.
  \item Les filtres transforment une valeur et se chaînent.
  \item \code{if}, \code{for}, \code{with}, \code{url} et \code{include} couvrent
        la majorité des gabarits.
  \item \code{csrf\_token} est obligatoire dans les formulaires \code{POST}.
  \item Les balises personnalisées vivent dans \code{templatetags/}.
  \item L'échappement automatique protège du XSS ; \code{safe} doit rester
        exceptionnel.
\end{itemize}
\end{retenir}
